\section*{Appendix A. Consolidated Stage 1 interview prompts (reference protocol)}
\label{app:stage1}

Across the 18 sessions, numbering and exact wording varied with pilot revisions; the following consolidated list reflects the most complete linear protocol used (e.g., Participant C) plus common follow-ups seen in other transcripts. Not every participant received every probe in identical order.

\paragraph{Musical taste and listening.}
How did you get into music? How did you get into metal, and what inspired you? Which metal genres and subgenres do you listen to, and why? Which other genres do you listen to besides metal?

\paragraph{Musical background and activity.}
Do you play instruments? Do you compose music? Do you have formal musical education? Are you currently in a band (name, tenure, role)? Have you been in other bands? How often do you perform live? How often do you record (studio or home studio)? Do you notice differences between live and studio vocal approach? How long have you been active as a singer or extreme vocalist?

\paragraph{Learning and technique.}
How did you learn extreme vocal techniques? Have you taken vocal lessons; did they cover extreme techniques? Do you follow warm-ups or routines? What strategies protect your voice from strain? Have you experienced vocal health problems related to extreme singing? How would you describe your vocal style? Which bands or vocalists influenced you? Which techniques do you use most frequently?

\paragraph{Perception, intelligibility, melody, and rhythm.}
Is intelligibility an important goal? How do you make lyrics intelligible with extreme techniques? Why might non-metal listeners find extreme vocals hard to understand? Why might some people dislike metal? How do you think about pitch, melody, and rhythm in extreme vocals? Can extreme vocals convey melody? Can they convey rhythm?

\paragraph{Additional probes used in some sessions (examples).}
Follow-up prompts such as subgenre detail, gateway artists, non-metal taste, gutturals and melody, combined melody/rhythm blocks, women in metal, historical change in bands, and open reflection on emotion and identity appeared when the conversation required clarification or depth.

\section*{Appendix B. Stage 2 (Qualtrics) items not represented in Appendix A}
\label{app:stage2}

The online survey added structured demographics, screening, inventories, and a dedicated perceptual-class block. Full consent language was shown per jurisdiction in Qualtrics; below are substantive research items only. Items Q47--Q53 in the survey overlap thematically with intelligibility/melody/rhythm prompts in Stage 1 and are omitted here to avoid duplication; they were retained in the dataset for cross-format comparison.

\paragraph{Ethics and screening.}
Origin-based consent routing; general participation consent; jurisdiction-specific consent where applicable; screening: ``Are you an active or former metal vocalist?''

\paragraph{Demographics.}
Age; gender (including free-text when ``other''); native language(s) learned at home; country of residence (optional).

\paragraph{Metal and non-metal taste (structured).}
Which metal genres and subgenres do you listen to regularly (multi-select, including other-specify)? Why do you like these genres? Which non-metal genres do you listen to (multi-select, including other-specify)? Why do you like those genres?

\paragraph{Musicianship and creative practice.}
Self-reported musicianship level; instruments and years of experience; whether you compose; description of composing workflow and tools; formal musical education (yes/no) and narrative description of training.

\paragraph{Band history and performance context.}
Current band membership and roles; tenure; past bands (names, roles, years); live performance frequency; recording experience in professional and/or home studios (narrative).

\paragraph{Vocal career, training, and health (survey form).}
Years active as extreme vocalist; formal vocal lessons (yes/no) and descriptions of private lessons, online learning, and peer learning; how extreme techniques were learned (structured choices plus other-specify); vocal exercise and health routines; strategies to prevent strain; history of vocal health problems (narrative); self-description of vocal style.

\paragraph{Influence and technique use (open and checklist).}
Free-text: which vocalists or bands influenced your style. Checklist: which labeled techniques you use most frequently (with other-specify); additional techniques in free text; differences between live and studio approach.

\paragraph{Perceptual-class focal block (Q43--Q46).}
Agreement with grouping harsh vocals into scream-like, growl-like, and hybrid/yell-like classes; optional comment on revising the system; matrix mapping each listed production technique to one or more classes (or not sure); open associations (emotions, imagery) for each class; self-reported frequency of thinking in those class terms while singing.

\paragraph{Historical and meta items (examples).}
Earliest recalled examples of growl, scream, and hybrid/yell vocals in metal or broader music history; whether animal comparisons help describe extreme vocals; open comments on classification and research; additional critiques or suggestions.
